\scanpage{51}

\unsure{corresponding|correspond} to systems of isomorphisms.
\[
\mu_{\mathcal D^*}:f_{\mathcal D^*}\longrightarrow f'_{\mathcal D^*},\qquad
\unsure{\mu_{U^*}}\;f_U\longrightarrow g_U.
\]
This corresponds \uline{conditions} $\tau'=\tau$, $d'_{i'}=d_{i'}$ $\forall i'\in I'$ and \unsure{to} the datum of
\inserted{this follows from the other conditions\unsure{\ldots}}
\margin{\unsure{let us express} the \uline{existence} of an isomorphism $\mu_{\mathcal D^*}$}

\margin{describes $\mu_U$}
a) $\boxed{\boxed{\mu\in\mathcal E}}$, defining an isomorphism between
\[
f_{\mathcal E}:\mathcal E'\longrightarrow\mathcal E
\quad\text{and}\quad
f'_{\mathcal E}:\mathcal E'\longrightarrow\mathcal E
\]
i.e. \deleted{\unsure{making}} such that
\[
\boxed{f'_{\mathcal E}=\operatorname{int}(\mu)f_{\mathcal E}}\,)
\]

b) $\boxed{\boxed{\mu_{i'}\in\mathcal E_i}}$ $(i'\in I',\ i=\tau(i'))$
\margin{describes $\mu_{\mathcal D^*}$}
satisfying
\[
\boxed{f'_{i'}=\operatorname{int}(\mu_{i'})f_{i'}}\,),
\]
\margin{this already implies $d_{i'}=d'_{i'}$}
the $\underset{\substack{\in\\\mathcal E}}{\mu},\ \underset{\substack{\in\\\mathcal E_i}}{\mu_{i'}}$ being related to the $\underset{\substack{\unsure{\in|\notin}\\\mathcal E}}{\alpha},\alpha'$ \unsure{\deleted{in}|in} $(\in\mathcal E)$
to the $\alpha_{i'},\alpha'_{i'}$ $(\in\Gamma)$ in the following way
\[
\boxed{\alpha'=p(\mu)\alpha}
\]
\[
\boxed{\alpha'_{i'}=\rho_i(\mu_{i'})\alpha_{i'}}\qquad\forall i'\in I'
\]
or again
\[
\boxed{\begin{cases}
p(\mu)=\alpha'\alpha^{-1}\\
\rho_i(\mu_{i'})=\alpha'_{i'}\alpha_{i'}^{-1}
\end{cases}}
\]
\margin{50}
\scanpage{52}
\mspage{24}

Note that, since $\rho_i$ is injective, the second relations, giving the $\rho_i(\mu_{i'})$, determine the $\mu_{i'}$ in a \uline{unique way}, the existence of these $\mu_{i'}$ is equivalent to the relations
\[
\boxed{\alpha'_{i'}\alpha_{i'}^{-1}\in\rho_i(\mathcal E_i)}\qquad i'\in I'.
\]

As for the relation $p(\mu)=\alpha'\alpha^{-1}$, it determines $\mu$ modulo multiplication \inserted{(on the right, say)} by a \deleted{\illegible{1}} $\mu_0\in\operatorname{Ker}p=\pi$, but which must be such that we still have
\[
f'_{\mathcal E}=\operatorname{int}(\deleted{\mu_0}\mu\mu_0)f_{\mathcal E}
\]
(in addition to $f'_{\mathcal E}=\operatorname{int}(\mu)f_{\mathcal E}$) which means \deleted{requires} that
\[
\operatorname{int}(\mu_0)f_{\mathcal E}=\unsure{f_{\mathcal E}|f'_{\mathcal E}}
\]
i.e.
\[
\boxed{\mu_0\in\operatorname{Centr}_{\unsure{\pi}}\deleted{\operatorname{Im}}f_{\mathcal E}(\mathcal E')}
\]

So that is the exact indeterminacy in the choice of a $0$-morphism \unsure{the|of the} $1$-morphisms \deleted{between} $f$ and $f'$ from $\mathcal D'$ to $\mathcal D$.
If for example we are under the conditions where
$\pi'\longrightarrow\pi$ induced by $f_{\mathcal E}:\mathcal E'\longrightarrow\mathcal E$ \deleted{is} has an image of finite index (open image)
--- case of a ``nonconstant morphism''! ---
then $\mu_0$ must centralize an open subgroup of
\margin{51}
\scanpage{53}

$\pi$ --- unless I am mistaken this too implies $\mu_0=1$, so it does seem that \deleted{\illegible{1}} (in the cases corresponding to dominant homomorphisms of algebraic curves) a \deleted{homomorphism} $0$-morphism from an $f$ \unsure{to} an $f'$ (if it exists) is unique, i.e. \uline{$\operatorname{Hom}$} $(\mathcal D',\mathcal D^*)$ is a \uline{discrete} category.

We can interpret $\mathcal D$ as $\mathcal E$ equipped with subgroups $\mathcal E_i\subset\mathcal E$, and with $p:\mathcal E\longrightarrow\unsure{\Gamma_*|\Gamma}$ (kernel $\pi$) finally with isomorphisms
\[
k_i:\mathrm T(\overline K)\xrightarrow{\sim}\mathcal E_i\cap\pi=L_i
\]
(commuting with the action of $\mathcal E_i$).
If we have \unsure{a|a} different system
\[
(\mathcal E',\ (\mathcal E'_{i'})_{i'\in I'},\ p':\mathcal E'\longrightarrow\Gamma',\ (k'_{i'})_{i'\in I'}),
\]
\inserted{$\mathcal E'_{i'}\subset\mathcal E'$}
a morphism from the first to the second is \uline{defined} by a
\[
f:\mathcal E'\longrightarrow\mathcal E
\]
satisfying the following conditions

\inserted{b) $pf$ is a conjugate of $f'$ [let $\alpha\in\Gamma$ be such that $pf=\operatorname{int}(\alpha)p'$]}

a) $\forall i'\in I',\quad\exists i\in I$ such that $f(\mathcal E'_{i'})$ is contained in a conjugate of $\mathcal E_i$
[let $\alpha_{i'}\in\mathcal E$ be such that
$f(\mathcal E'_{i'})\unsure{\subset|=}\operatorname{int}(\alpha_{i'}^{-1})(\mathcal E_i)$]
\inserted{i.e. $\operatorname{int}(\alpha_{i'})f(\mathcal E'_{i'})\subset\mathcal E_i$}
\margin{52}
\scanpage{54}
\mspage{25}

\uline{N.B.} The $i\in I$ corresponding to $i'\in I'$ is unique, whence $\tau:I'\longrightarrow I$.
If the $\alpha_{i'}$ are chosen, we deduce \deleted{\unsure{$\forall i'$}} induced homomorphisms
$f_{i'}:\mathcal E'_{i'}\longrightarrow\mathcal E_i$,
\[
f_{i'}=(\operatorname{int}(\alpha_{i'})\circ f)|\mathcal E'_{i'},
\]
which then induces, via the $k'_{i'},k_i$, an endomorphism
\[
f^0_{i'}:\mathrm T(\overline K)\longrightarrow\mathrm T(\overline K).
\]
\margin{\uline{N.B.} A priori $f^0_{i'}$ is \unsure{multiplication by} $\chi_{i'}\in\unsure{\widehat{\mathbb Z}^{\,*}|\widehat{\mathbb Z}}$. The centralizer in $\Gamma$ of an open subgroup is reduced to the identity.}
We require that $\alpha\in\Gamma$ (which conjugates $\unsure{f_{\mathcal E}}$ to $\unsure{f'_{\mathcal E}}$\}, and is \unsure{probably} uniquely determined by this condition) and $\alpha_{i'}$ (which is probably determined modulo \deleted{\unsure{conjugation}} composition on the left with an element \unsure{of the normalizers|of the normalizer} of $\mathcal E_i$ (at least of the transporters \inserted{of the $\mathcal E_i$} of an open subgroup of $\mathcal E_i$ --- it is the same if $f$ is an isomorphism) can be chosen in such a way that, putting
\[
\beta_{i'}=p(\alpha_{i'})\alpha,
\]
we have
\[
f^0_{i'}=(d_{i'}\chi(\beta_{i'}))\operatorname{id}_{\mathrm T(\overline K)}
\]
\inserted{i.e. $\chi_{i'}=d_{i'}\chi(\beta_{i'})$}
\margin{This condition on the systems of the $\alpha_{i'}$ does not change ($\alpha$ remaining fixed) if we change $\alpha_{i'}$ to $\mu_{i'}\alpha_{i'}$, where $\mu_{i'}\in\mathcal E_i$.}
$d_{i'}\in\mathbb N^*$ is a natural integer obviously uniquely determined \unsure{\}} when $\alpha$ and the $\alpha_{i'}$ are chosen).

\deleted{Let us replace $\alpha$ by $\alpha'=\alpha_0\alpha$.}
\margin{53}
\scanpage{55}

8, \uline{Taxonomic reflection.}

The case $\dim X=0$ translates in the paradigm of the system
$(\mathrm T(\overline K)\xrightarrow{k_i}\mathcal E_i\overset{\rho_i}{\hookrightarrow}\mathcal E\xrightarrow{p}\Gamma)$
\deleted{\illegible{2}} into the condition \deleted{$\unsure{\pi|I}=\varnothing$} $p$ injective, \deleted{thus \unsure{is|of}}
\inserted{and thus $I=\varnothing$ --- so this case is}
described \deleted{\illegible{1}}
simply by giving an open \uline{subgroup} $\mathcal E\subset\Gamma$, i.e. a finite subextension $\unsure{L}$ of $\overline K/K$
(N.B. of course we will have $X'=\operatorname{Spec}\unsure{L'|L}$, and the choices \inserted{made} on $X$ \inserted{--- i.e. on $\underset{=\varnothing}{\Pi_{\mathcal D^*}}\longrightarrow\Pi_U\longrightarrow\Pi_e$ ---} \unsure{leading|lead} to this description --- amount here to the choice of such a \inserted{$K$-}embedding of $L\deleted{\unsure{K'}}=\mathrm H^0(X,\mathcal O_X)$ into $\overline K/K$)

The case $\dim X=1$ translates into the fact that $\mathcal E\longrightarrow\Gamma$ is not injective --- as for $I$, it may be empty or not in this case.
If it is empty, the description is thus simply in terms of a homomorphism of profinite groups $\mathcal E\xrightarrow{p}\Gamma$ (where $\Gamma$ given in advance $=\operatorname{Gal}(\overline K/K)$ \unsure{\}} $=\operatorname{Aut}_{\Pi_e}(\Omega)$).

Suppose now given both
\[
\mathcal D=(\mathcal E,p,I,(\rho_i:\mathcal E_i\hookrightarrow\mathcal E),(k_i:\mathrm T(\overline K)\longrightarrow\mathcal E_i))
\]
and
\[
\mathcal D'\unsure{=} (\mathcal E',\unsure{p'},I',(\rho'_{i'}:\ ),(k'_{i'}:\ )),
\]
and let us return to the question of the description of the morphisms from $\mathcal D'$ to $\mathcal D$.
We are going to distinguish four cases, according to the two values
\margin{54}
\scanpage{56}
\mspage{26}

$0,1$ possible for $n=\dim\mathcal D$ and $n'=\dim\mathcal D'$ respectively.

I) $n=0,n'=0$, \deleted{\unsure{then|thus}} $I=I'=\varnothing$ and the data reduce to open subgroups
$\mathcal E\overset{p}{\hookrightarrow}\Gamma$, $\mathcal E'\overset{p'}{\hookrightarrow}\Gamma$.
In the cumbersome description above, \deleted{all} the considerations relating to $I,I'$ drop out and \uline{the}\margin{(a)} morphism $\mathcal D'\longrightarrow\mathcal D$ amounts to a class of \deleted{\unsure{homomorphisms}} pairs $(f,\alpha)$, where
$f:\mathcal E'\longrightarrow\mathcal E$ and $\alpha\in\Gamma$, such that we have
\[
pf=\operatorname{int}(\alpha)f
\]
--- which implies that $f$ is already determined by $\alpha$, being induced by $\operatorname{int}(\alpha)$
\inserted{[$\alpha$ subject to $\alpha\in\operatorname{Trans}_{\Gamma}(\mathcal E',\mathcal E)$ \unsure{$=$} $\operatorname{Trans}_{\Gamma}(L,L')$]}
\deleted{\diagram{\begin{tikzcd}[ampersand replacement=\&,column sep=1.5em,row sep=2em]
\unsure{\mathcal E'}\arrow[rr,"f"]\arrow[dr,"p'"'] \& \& \mathcal E\arrow[dl,"p"] \\
\& \Gamma \&
\end{tikzcd}}}
\diagram{\begin{tikzcd}[ampersand replacement=\&]
\mathcal E' \arrow[r,"f"] \arrow[d,"p'"'] \& \mathcal E \arrow[d,"p"] \\
\Gamma \arrow[r,"\operatorname{int}(\alpha)","\sim"'] \& \Gamma
\end{tikzcd}}
[which induces \unsure{thus|indeed} a homomorphism from $\unsure{L^*|L}\subset\overline K$ to $L'\subset\overline K$, hence a homomorphism $\operatorname{Spec}(L')\longrightarrow\operatorname{Spec}(L)$ in the opposite direction].

The condition that $\alpha,\alpha'\in\Gamma$ define the same morphism $\mathcal D'\longrightarrow\mathcal D$ (or again $\operatorname{Spec}\unsure{L'}\longrightarrow\operatorname{Spec}L$) is described by the existence of a $\mu\in\mathcal E$ \deleted{(= stabilizer of} such that $\alpha'=\mu\alpha$.
The fact that we give a $1$-$1$ correspondence between
\[
\operatorname{Hom}_{K}(X',X)\simeq\unsure{\operatorname{Hom}}_{K}(\unsure{L},L')
\]
with $\operatorname{Hom}(\mathcal D',\mathcal D)$ made explicit too, is clear.
\margin{55}
\scanpage{57}

II) $n'=1,n=0$.
\inserted{Here $n=0$ already implies $I=\varnothing$.}
\deleted{Here} For $\operatorname{Hom}(X',X)\ne\varnothing$, we must have $I'=\varnothing$
[by the condition $f^{-1}(S)=S'$ which implies $S'=\varnothing$, if $f:(X',S')\longrightarrow(X,S)$]
\inserted{$f^{-1}(\varnothing)=\varnothing$}
--- \deleted{in} in the description of the morphisms $\mathcal D'\longrightarrow\mathcal D$, this corresponds to the fact that we can have a map $\tau:I'\longrightarrow I(=\varnothing)$ only if $I'=\varnothing$.
\margin{this we cannot}
\deleted{\illegible{1}}
Restricting ourselves therefore to the case $I'=\varnothing$ (otherwise $\operatorname{Hom}(X',X)=\operatorname{Hom}(\mathcal D',\mathcal D)=\varnothing$ and we are happy), we thus find that $\mathcal D$ amounts to giving $\mathcal E\overset{p}{\hookrightarrow}\Gamma$ an open subgroup, and $\mathcal D'$ to giving a homomorphism of profinite groups $p':\mathcal E'\longrightarrow\Gamma$, whose image will be an open subgroup of $\Gamma$, say $\Gamma'_0\subset\Gamma$, corresponding to a \inserted{finite} subextension $L'$ of $\overline K/K$.
\margin{Restricting ourselves therefore to the case}

The choices made concerning $X'$ imply \unsure{precisely|that we have} a fixed isomorphism $L'\simeq\mathrm H^0(X',\mathcal O_{X'})$.
That said, a $K$-morphism \inserted{from $(X',\varnothing)$ to $(X,\varnothing)$ i.e. from} $X'\longrightarrow X$ amounts \inserted{(since $X$ affine)} to giving a $K$-morphism $\operatorname{Spec}L'\longrightarrow X\ (\simeq\operatorname{Spec}L)$ from the affine envelope, i.e. a \inserted{$K$-}homomorphism $L\longrightarrow L'$.
(which depends only on $L,L'$, or again \unsure{on the} \deleted{\illegible{1}} subgroups $\mathcal E\subset\Gamma$ and $\Gamma_0\subset\Gamma$
\margin{56}
\scanpage{58}
\mspage{27}

of $\Gamma$).
On the other hand the description in terms of the diagrams $\mathcal D,\mathcal D'$ amounts to saying that a morphism is determined by a pair $(f,\alpha)$, $f:\mathcal E'\longrightarrow\mathcal E$ and $\alpha\in\Gamma$, such that we again have commutativity
\diagram{\begin{tikzcd}[ampersand replacement=\&]
\mathcal E' \arrow[r,"f"] \arrow[d,"p'"'] \& \mathcal E \arrow[d,hook,"p"] \\
\Gamma \arrow[r,"\operatorname{int}(\alpha)"] \& \Gamma
\end{tikzcd}}
\[
pf=\operatorname{int}(\alpha)p'
\]
which again implies that $f$ is determined by $\alpha$ (namely, induced by $\operatorname{int}(\alpha)$), and $\alpha$ being subject to the sole condition
\[
\alpha\in\operatorname{Transp}_{\Gamma}(\Gamma'_0,\mathcal E)
\]
--- and $\alpha,\alpha'$ \unsure{describing|describe} the $=$ homomorphism \unsure{iff}
\[
\exists\mu\in\mathcal E\quad\text{such that}\qquad\unsure{\Longleftarrow|\Longleftrightarrow}
\]
\[
\alpha'=\mu\alpha
\]
so we again have
\[
\begin{aligned}
\operatorname{Hom}(\mathcal D',\mathcal D)
&\simeq\mathcal E\backslash\operatorname{Transp}_{\Gamma}(\Gamma'_0,\Gamma)\ \deleted{\illegible{1}}\\
&\simeq\operatorname{Transp}_{\Gamma}(L,L')/\operatorname{Fix}_{\Gamma}(L)\\
&\simeq\operatorname{Hom}_{\unsure{K\text{-}\mathrm{ext}|K\text{-}\mathrm{cat}}}(L,L')
\end{aligned}
\]
so we again find in this case \unsure{the|the} full faithfulness.

In fact, we could have treated together
\margin{57}
\scanpage{59}

the two cases $(n'=0,n'=1)$ where $n=0$, and the description of the $\operatorname{Hom}(X,X')$ \inserted{in terms of $\mathcal D,\mathcal D'$} is valid moreover, without \deleted{assuming that} making draconian hypotheses (projective, smooth, $\dim\leq1$, anabelian) on $X'$.
\deleted{It suffices that $X'$ be a connected quasi-compact quasi-separated scheme over $K$, \unsure{thus|in} the morphisms $X'\longrightarrow X=\operatorname{Spec}(L)$ correspond to the morphisms $\operatorname{Spec}A'\longrightarrow X$ from the affine envelope of $X'$ to $X$, defined as the Spectrum of $A'=\mathrm H^0(X',\mathcal O_{X'})$, i.e. to the \unsure{$K$-homomorphisms} of $K$-algebras
\[
L\longrightarrow A'.
\]}

Indeed, since $X$ is \inserted{finite} étale over $K$, its inverse image \inserted{$X_{X'}=$} \deleted{\illegible{1}} $X'\times_K\unsure{X}$ over $X'$ is \inserted{finite} étale over $X'$, and
\[
\operatorname{Hom}_K(X,X')\simeq\deleted{\operatorname{Hom}\unsure{(}}\operatorname{Sections}(X_{X'}/X')
\]
can be described in terms of the category $\widehat{\Pi_1(X')}$ $(=\Pi_U)$ of finite local systems on $X'$ i.e. of finite local systems on $\Pi_U$.
Following this description, we arrive at the
\margin{58}
\scanpage{60}
\mspage{28}

common description given in I, II, without any other hypothesis on $X$ than \deleted{connectedness} $0$-connectedness (to be able to describe $\Pi_U=\Pi_{X'}$ over $\Pi_e$ in terms of a group homomorphism $p':\mathcal E'\longrightarrow\Gamma\ldots$).

When we are given $(X',S')$, with $S'$ not necessarily empty, $X'$ $0$-connected (and \deleted{dimension} $X'$ smooth of dimension $1$ at the points of $S'$ (it would suffice $=$ that $S'$ \unsure{not} \unsure{disconnect|disconnect} locally \inserted{(étale)} $X'$, $\ldots$),
\inserted{thus $X'-S'$ $0$-connected}
then the \deleted{part} datum
$\Pi_1(X'-S')\longrightarrow\Pi_1(e_K)$ is described \deleted{\illegible{1}} by $\mathcal E'\xrightarrow{p'}\Gamma$ (independently of the consideration of the $\mathcal E_i$ etc) and hence $\operatorname{Hom}_K(X'-S',X)$ is described as we just said.
But we have (with the normality hypotheses on $X'$ at the points of $S'$)
\[
\operatorname{Hom}_K(X'-S',X)=\operatorname{Hom}_K(X',X).
\]

This encourages modifying the definition of morphisms $(X',S')\longrightarrow(X,S)$ given \unsure{at the|by} beginning, via $f:X'\longrightarrow X$ with $f^{-1}(S)_{\mathrm{r\acute ed}}=S'$, by requiring this \inserted{not for $f$, but} \unsure{only} for the restriction of $f$ \unsure{to} the union of the irreducible components of $X'$
\margin{59}
